% !TeX root = ../applied-mathematics1-companion.tex
\chapter{Continuous Mathematical Models}\label{chap:applied-continuous}

\section{From a verbal model to an ODE}
Continuous time describes a state through a rate of change:
\[
\dot x=\frac{dx}{dt}=F(x,t),\qquad x(0)=x_0.
\]
An equation is \emph{autonomous} when its right-hand side has no explicit
time dependence: $\dot x=f(x)$. Distinguish this from the discrete equation
$x_{t+1}=f(x_t)$: the same expression on the right can give very different
behaviour in the two settings.

\studyguide{Choose the state variables and units; state assumptions;
identify gains and losses; write rate of change = gains minus losses;
check units and physical constraints; impose initial conditions; then solve
or analyse equilibria and signs.}
For a scalar state of units $U$, the right side of an ODE must have units
$U/\text{time}$. This simple check often catches a missing factor.
The models here have smooth right sides, so standard local existence and
uniqueness results apply. A solution is followed only while it exists in
the relevant domain; not every nonlinear model has solutions for all time.

\subsection{Partial derivatives and first-order expansion}
For $f(x,y)$, partial differentiation varies one argument while fixing the
other. In symbols,
\[
f_x(x,y)=\lim_{h\to0}\frac{f(x+h,y)-f(x,y)}h,
\quad
f_y(x,y)=\lim_{h\to0}\frac{f(x,y+h)-f(x,y)}h.
\]
For instance, $f=x^2y+\cos y$ gives $f_x=2xy$ and $f_y=x^2-\sin y$.
For a differentiable function,
\[
f(a+X,b+Y)=f(a,b)+f_x(a,b)X+f_y(a,b)Y
+o\bigl(\sqrt{X^2+Y^2}\bigr).
\]
The derivatives evaluated at $(a,b)$ are constants in the linear approximation.

\section{Scalar population models}
\subsection{Exponential growth}
If growth rate is proportional to the population, $\dot P=kP$, $k>0$.
The solution is $P(t)=P_0e^{kt}$, including $P\equiv0$ when $P_0=0$.
Positive initial populations increase indefinitely. This can describe a
limited time range but ignores eventual resource constraints.

\subsection{Continuous logistic growth}
Competition within a population suggests
\begin{equation}\label{eq:applied-cont-logistic}
\dot P=kP\left(1-\frac PK\right),\qquad k,K>0.
\end{equation}
Here $k$ has units inverse time and $K$ has population units. The equilibria
are $0$ and $K$. For $0<P<K$, $\dot P>0$; for $P>K$, $\dot P<0$.
Thus the phase line points towards $K$ on both sides in the positive domain.

To solve, first treat the equilibria separately, since separation divides
by $P(K-P)$. For other initial values,
\begin{align*}
\frac{dP}{P(K-P)}&=\frac{k}{K}\,dt,\\
\frac1K\left(\frac1P+\frac1{K-P}\right)dP&=\frac{k}{K}\,dt,\\
\log\left|\frac P{K-P}\right|
-\log\left|\frac{P_0}{K-P_0}\right|&=kt.
\end{align*}
The sign of $P/(K-P)$ is preserved along a non-equilibrium trajectory.
After applying the initial condition and rearranging,
\begin{equation}\label{eq:applied-logistic-solution}
P(t)=\frac{KP_0}{P_0+(K-P_0)e^{-kt}}.
\end{equation}
This formula also gives both equilibrium solutions when $P_0=0$ or $K$.
For $P_0>0$, its denominator remains positive for $t\geq0$ and
$P(t)\to K$. If $P_0>K$, absolute values in the logarithmic derivation
are essential, even though the final formula has the same form.

\section{Phase lines and scalar stability}
For $\dot P=f(P)$, an equilibrium satisfies $f(P^*)=0$.
Write $P=P^*+p$ and expand:
\[
\dot p=f(P^*+p)=f'(P^*)p+o(p).
\]
The linearised equation $\dot p=\lambda p$ has solution
$p(t)=p_0e^{\lambda t}$, with $\lambda=f'(P^*)$.
For a continuously differentiable $f$ near $P^*$:
\begin{itemize}
\item $f'(P^*)<0$ implies local asymptotic stability;
\item $f'(P^*)>0$ implies instability;
\item $f'(P^*)=0$ leaves the linear test inconclusive.
\end{itemize}
The condition here is a \emph{negative derivative}, not derivative modulus
less than $1$. The latter belongs to discrete maps.

\begin{example}[A periodic array of equilibria: source Example 2.1]
For $\dot P=\sin P$, equilibria satisfy $P^*=n\pi$, $n\in\mathbb Z$.
Since $f'(P^*)=\cos(n\pi)=(-1)^n$, odd multiples attract and even multiples
are unstable. Between neighbouring even multiples $2j\pi$ and
$2(j+1)\pi$, every non-equilibrium solution tends to $(2j+1)\pi$.
If $P$ denotes a population, restrict the mathematical phase line to
nonnegative values.
\end{example}

\begin{example}[Additional handwritten example, 15 October]
Consider $\dot x=f(x)=e^{-x}\sin x$ for initial values $x_0\in[0,6]$.
The equilibria in this interval are $0$ and $\pi$.
Since $f'(x)=e^{-x}(\cos x-\sin x)$,
\[
f'(0)=1>0,\qquad f'(\pi)=-e^{-\pi}<0.
\]
Thus $0$ is unstable, while $\pi$ attracts. For $0<x<\pi$ the state moves
right; for $\pi<x\leq6<2\pi$ it moves left. Every $x_0\in(0,6]$ tends to
$\pi$, while $x_0=0$ stays fixed.

To interpret the changing speed, note $|\dot x|=|f(x)|$.
The positive velocity has its maximum at $x=\pi/4$; the negative velocity
has its minimum at $x=5\pi/4$. A right-moving state below $\pi/4$ first
speeds up and then slows down. A left-moving state above $5\pi/4$ first
speeds up as it approaches $5\pi/4$, then slows towards $\pi$.
A negative velocity and a decreasing speed are different statements.
\end{example}

\section{A mixing model}
A tank initially contains $100$ litres of pure water. Brine with concentration
$5$ grams per litre enters at $2$ litres per minute; an equal outflow keeps
the volume constant. Assume perfect instantaneous mixing.
Let $M(t)$ be the salt mass in grams and $t$ be measured in minutes.
\[
\text{inflow salt rate}=5\times2=10\ \mathrm{g/min},\qquad
\text{outflow salt rate}=\frac{M}{100}\times2=\frac{M}{50}\ \mathrm{g/min}.
\]
Therefore $\dot M=10-M/50$, with $M(0)=0$. An integrating factor gives
\begin{align*}
\frac d{dt}\left(Me^{t/50}\right)&=10e^{t/50},\\
M(t)&=500+Ce^{-t/50},\\
M(0)=0&\quad\Longrightarrow\quad C=-500.
\end{align*}
The mass and concentration are
\[
M(t)=500(1-e^{-t/50})\ \mathrm g,\qquad
c(t)=\frac{M(t)}{100}=5(1-e^{-t/50})\ \mathrm{g/L}.
\]
As $t\to\infty$, the tank concentration approaches the incoming concentration.
If inflow and outflow volumes differed, the volume would vary and the
outflow concentration would be $M(t)/V(t)$ instead.

\section{Coupled systems, phase planes and nullclines}
A two-species model has the form
\[
\dot x=f(x,y),\qquad\dot y=g(x,y).
\]
Terms proportional to $xy$ represent encounters under a well-mixed
assumption. Their signs encode the interaction. For positive constants:
\begin{center}
\begin{tabular}{lll}
\toprule
Interaction&$\dot x$&$\dot y$\\\midrule
Predation ($x$ prey, $y$ predator)&$ax-bxy$&$-cy+dxy$\\
Competition&$ax-bxy$&$cy-dxy$\\
Mutual benefit&$ax+bxy$&$cy+dxy$\\\bottomrule
\end{tabular}
\end{center}
These are simplified models; sign information alone does not capture all
resource limitations or guarantee bounded populations.

\begin{definition}[Phase portrait and nullclines]
A trajectory is the curve traced by $(x(t),y(t))$ in the phase plane,
with arrows showing increasing time. The vector field is $(f(x,y),g(x,y))$.
The $x$-nullclines satisfy $f(x,y)=0$, where nonzero arrows are vertical.
The $y$-nullclines satisfy $g(x,y)=0$, where nonzero arrows are horizontal.
Equilibria satisfy both equations simultaneously.
\end{definition}
Under uniqueness, two distinct trajectories cannot cross at a non-equilibrium
point: the same state cannot have two different future solutions.
A nullcline is generally not a trajectory, and a point on just one
nullcline need not be an equilibrium.

\begin{example}[Reading the source phase portrait: source Example 2.3]
The supplied sketch has rightward flow on the $x$-axis and downward flow
on the $y$-axis, with an unstable origin, an unstable interior equilibrium,
and a closed orbit attracting nearby trajectories. It illustrates an
attracting \emph{limit cycle}: an isolated periodic orbit.
One cannot infer this entire portrait from interaction signs alone.
In particular, the closed Lotka--Volterra orbits studied below are not
attracting limit cycles, because they form a family of neighbouring cycles.
\end{example}

\studyguide{For a phase portrait: solve $f=g=0$; draw both types of nullcline;
check signs of $f$ and $g$ in each region; combine them into flow arrows;
then draw trajectories consistent with those directions and uniqueness.}

\begin{example}[Competing species: source Section 2.4.2]\label{ex:applied-competition}
Consider
\[
\dot x=x(5-x-y),\qquad\dot y=y(8-x-2y).
\]
The $x$-nullclines are $x=0$ and $x+y=5$; the $y$-nullclines are $y=0$
and $x+2y=8$. Their intersections give
\[
(0,0),\qquad(5,0),\qquad(0,4),\qquad(2,3).
\]
For $x,y>0$, $\dot x$ has the sign of $5-x-y$, and $\dot y$ has the sign
of $8-x-2y$. At $(1,1)$, both are positive, so arrows point up and right.
Above both oblique nullclines, both are negative, so arrows point down and left.
The axes are invariant because the corresponding population derivative vanishes.
\end{example}
% !TeX root = ../applied-mathematics1-companion.tex
\begin{center}
\begin{tikzpicture}[x=.8cm,y=.8cm,>=Stealth]
\draw[->] (0,0)--(8.8,0) node[right] {$x$};
\draw[->] (0,0)--(0,5.8) node[above] {$y$};
\foreach \x in {.5,1.5,2.5,3.5,4.5,5.5,6.5,7.5}
 \foreach \y in {.5,1.5,2.5,3.5,4.5} {
 \pgfmathsetmacro{\dx}{\x*(5-\x-\y)}
 \pgfmathsetmacro{\dy}{\y*(8-\x-2*\y)}
 \pgfmathsetmacro{\norm}{max(sqrt(\dx*\dx+\dy*\dy),.01)}
 \draw[->,muted!60] (\x,\y)--++({.24*\dx/\norm},{.24*\dy/\norm});}
\draw[thick,ink] (0,5)--(5,0);
\draw[thick,accent,dashed] (0,4)--(8,0);
\node[ink,above right] at (.2,4.8) {$x+y=5$};
\node[accent,above right] at (5,1.5) {$x+2y=8$};
\foreach \p in {(0,0),(5,0),(0,4),(2,3)} \fill \p circle (2pt);
\node[below left] at (0,0) {$0$};
\node[below] at (5,0) {$5$};
\node[left] at (0,4) {$4$};
\node[above right,fill=white,inner sep=1pt] at (2,3) {$(2,3)$};
\end{tikzpicture}

{\small Competition model: oblique nullclines and normalised flow arrows.}
\end{center}


\section{Linearisation of two-dimensional systems}
At an equilibrium $(x^*,y^*)$, write $x=x^*+X$ and $y=y^*+Y$.
First-order expansion gives the linearised system
\[
\begin{pmatrix}\dot X\\\dot Y\end{pmatrix}
=J\begin{pmatrix}X\\Y\end{pmatrix},\qquad
J=\begin{pmatrix}A&B\\C&D\end{pmatrix}
=\left.\begin{pmatrix}f_x&f_y\\g_x&g_y\end{pmatrix}\right|_{(x^*,y^*)}.
\]
The matrix $J$ is the \emph{Jacobian at the equilibrium}. For the original
nonlinear system there are also higher-order remainder terms.

Try a mode $(X,Y)^{\mathsf T}=\mathbf v e^{\lambda t}$, with
$\mathbf v\ne\mathbf0$. Substitution gives $J\mathbf v=\lambda\mathbf v$.
Thus $\lambda$ is an eigenvalue and must satisfy
\begin{equation}\label{eq:applied-characteristic}
\det(\lambda I-J)=(\lambda-A)(\lambda-D)-BC=0.
\end{equation}
This determinant approach works even when $B=0$ or $C=0$; dividing by a
matrix entry could lose those cases.

For distinct eigenvalues with independent eigenvectors, the general linear
solution is
\[
\begin{pmatrix}X\\Y\end{pmatrix}
=c_1\mathbf v_1e^{\lambda_1t}+c_2\mathbf v_2e^{\lambda_2t}.
\]
There are two independent constants, not four unrelated coefficients in
$X$ and $Y$. Repeated eigenvalues may require a term involving $te^{\lambda t}$.
Complex conjugate modes combine into real sine and cosine terms.
Superposition applies to the homogeneous linearised system, not generally
to the original nonlinear model.

\begin{proposition}[Local linear stability]
For a continuously differentiable nonlinear vector field, if all eigenvalues
of its equilibrium Jacobian have negative real part, the equilibrium is
locally asymptotically stable. If at least one has positive real part,
it is unstable. If some real parts are zero and none are positive,
linearisation alone generally does not decide nonlinear stability.
\end{proposition}
These standard linearisation results justify using the modes as a local
stability test. The linear solution is not an exact formula for the full
nonlinear trajectory, even when it predicts stability correctly.

\subsection{Trace and determinant: a supplementary shortcut}
Let $\tau=A+D$ and $\Delta=AD-BC$. The characteristic equation is
$\lambda^2-\tau\lambda+\Delta=0$.
\begin{itemize}
\item $\Delta<0$: one positive and one negative eigenvalue, an unstable saddle.
\item $\Delta>0$, $\tau<0$: negative real parts, local attraction.
\item $\Delta>0$, $\tau>0$: positive real parts, instability.
\item $\Delta>0$, $\tau=0$: purely imaginary eigenvalues; no nonlinear
stability conclusion from the linear test alone.
\end{itemize}
For $\Delta>0$, the sign of $\tau^2-4\Delta$ distinguishes two real modes
from complex conjugate modes. The cases $\Delta=0$ require further analysis.

\subsection{Return to the competition example}
At $(2,3)$, direct substitution gives
\begin{align*}
\dot X&=-2X-2Y-X^2-XY,\\
\dot Y&=-3X-6Y-XY-2Y^2.
\end{align*}
Dropping the quadratic terms gives
\[
J=\begin{pmatrix}-2&-2\\-3&-6\end{pmatrix},\qquad
\lambda^2+8\lambda+6=0,\qquad\lambda=-4\pm\sqrt{10}.
\]
Both eigenvalues are negative, so coexistence is locally asymptotically stable.
For reference, the general Jacobian is
\[
J(x,y)=\begin{pmatrix}5-2x-y&-x\\-y&8-x-4y\end{pmatrix}.
\]
The other equilibria have eigenvalues $5,8$ at $(0,0)$; $-5,3$ at $(5,0)$;
and $1,-8$ at $(0,4)$. Thus the origin is repelling and the two nonzero
axis equilibria are saddles. These extra classifications complete the
local picture without asserting global attraction for every initial state.

\section{The Lotka--Volterra predator--prey model}
\textbf{Convention for this section: $x$ is predator population, $y$ is prey.}
This reverses the labels in the earlier illustrative interaction table.
With $a,b,c,d>0$,
\begin{equation}\label{eq:applied-lv}
\dot x=x(-a+by),\qquad\dot y=y(c-dx).
\end{equation}
Predators decline without prey, prey grow without predators, and encounters
benefit predators while reducing prey. Positive populations stay positive:
for example $x(t)=x_0\exp\!\left(\int_0^t[-a+by(s)]\,ds\right)$.
The axes are invariant and are analysed separately from the positive interior.

The nullclines are $x=0$, $y=a/b$ for $\dot x=0$, and $y=0$, $x=c/d$
for $\dot y=0$. Equilibria are $(0,0)$ and $(c/d,a/b)$.
Above $y=a/b$ predators grow; to the right of $x=c/d$ prey decrease.
For this predator-on-horizontal-axis convention, the positive orbits
run clockwise around coexistence.

\subsection{Harvesting}
Per-capita harvesting rates $h_1,h_2\geq0$ give
\[
\dot x=x(-a-h_1+by),\qquad\dot y=y(c-h_2-dx).
\]
If $c>h_2$, the positive equilibrium shifts to
\[
\left(\frac{c-h_2}{d},\frac{a+h_1}{b}\right).
\]
Thus its predator value falls and prey value rises relative to the original
coexistence equilibrium. These are equilibrium or cycle-average comparisons,
not a claim that every instantaneous prey population rises. If $c\leq h_2$,
there is no positive coexistence equilibrium of this form.

\subsection{Local analysis}
At the origin, the linearised equations are $\dot X=-aX$, $\dot Y=cY$,
so the origin is a saddle.
At coexistence,
\[
J=\begin{pmatrix}0&bc/d\\-ad/b&0\end{pmatrix},\qquad
\lambda^2=-ac,\qquad\lambda=\pm i\sqrt{ac}.
\]
Let $\omega=\sqrt{ac}$ and $B=bc/d$. The real linear solution can be written
\begin{align*}
X(t)&=A_1\cos(\omega t)+A_2\sin(\omega t),\\
Y(t)&=\frac\omega B\bigl[-A_1\sin(\omega t)+A_2\cos(\omega t)\bigr].
\end{align*}
It oscillates with period $2\pi/\sqrt{ac}$, using two independent constants.
This is the small-amplitude linearised period, not the exact period of
every nonlinear orbit. Purely imaginary eigenvalues alone do not prove
nonlinear stability; we next use a conserved quantity.

\subsection{A first integral and closed trajectories}
For $x,y>0$ and away from a vertical trajectory tangent, take the ratio:
\[
\frac{dy}{dx}=\frac{y(c-dx)}{x(-a+by)}.
\]
Separation yields
\[
\int\left(-\frac ay+b\right)dy
=\int\left(\frac cx-d\right)dx,
\]
so a trajectory satisfies
\begin{equation}\label{eq:applied-lv-integral}
c\log x-dx+a\log y-by=\text{constant}.
\end{equation}
Logs require \emph{strictly positive} populations, not merely nonnegative ones.
Direct differentiation verifies the conserved expression along the whole
positive trajectory, including points where the ratio $dy/dx$ was unavailable.
Exponentiating gives
\[
F(x)G(y)=H>0,\qquad F(x)=x^ce^{-dx},\quad G(y)=y^ae^{-by}.
\]
Each factor tends to zero at both ends of its positive domain. Also
\[
F'(x)=x^{c-1}e^{-dx}(c-dx),\qquad
G'(y)=y^{a-1}e^{-by}(a-by).
\]
Thus $F$ has a unique maximum $M_1=(c/d)^ce^{-c}$ at $c/d$;
$G$ has a unique maximum $M_2=(a/b)^ae^{-a}$ at $a/b$.
At $H=M_1M_2$, the level set is just the equilibrium. For
$0<H<M_1M_2$, horizontal and vertical lines through the equilibrium each
meet the level set twice, and the trajectory is a closed curve surrounding it.

\textit{Supplementary justification.} Equivalently,
\[
V(x,y)=dx-c\log x+by-a\log y
\]
is strictly convex, tends to $+\infty$ at the boundary and at infinity,
and has its unique minimum at coexistence. Its non-minimal level sets
are compact smooth closed curves with no equilibrium on them.
The nonvanishing vector field follows each such curve periodically.
It also follows that coexistence is stable in the non-attracting sense:
nearby level curves remain nearby, but no other one converges to the centre.
It is therefore a nonlinear \emph{centre}, not an asymptotically stable state
and not an attracting limit cycle.
% !TeX root = ../applied-mathematics1-companion.tex
% Newly integrated illustrative trajectories for a=b=c=d=1.
\begin{center}
\begin{tikzpicture}[x=1.8cm,y=1.8cm,>=Stealth]
\draw[->] (0,0)--(3.1,0) node[right] {$x$ (predators)};
\draw[->] (0,0)--(0,3.1) node[above] {$y$ (prey)};
\draw[muted,dashed] (1,0)--(1,2.9);
\draw[muted,dashed] (0,1)--(2.9,1);
\draw[ink,thick] plot[smooth] coordinates {(1.3000,1.0000) (1.3000,0.9988) (1.2995,0.9845) (1.2981,0.9705) (1.2958,0.9568) (1.2927,0.9433) (1.2888,0.9303) (1.2841,0.9176) (1.2786,0.9053) (1.2725,0.8934) (1.2656,0.8819) (1.2581,0.8709) (1.2500,0.8603) (1.2414,0.8502) (1.2322,0.8406) (1.2225,0.8315) (1.2124,0.8229) (1.2019,0.8147) (1.1911,0.8071) (1.1799,0.7999) (1.1684,0.7932) (1.1567,0.7871) (1.1448,0.7814) (1.1327,0.7762) (1.1205,0.7715) (1.1081,0.7673) (1.0957,0.7635) (1.0833,0.7603) (1.0708,0.7575) (1.0583,0.7551) (1.0459,0.7532) (1.0336,0.7518) (1.0213,0.7508) (1.0092,0.7503) (0.9971,0.7502) (0.9852,0.7505) (0.9735,0.7512) (0.9620,0.7524) (0.9507,0.7540) (0.9395,0.7559) (0.9287,0.7583) (0.9180,0.7611) (0.9076,0.7643) (0.8975,0.7679) (0.8876,0.7719) (0.8780,0.7762) (0.8688,0.7810) (0.8598,0.7861) (0.8511,0.7916) (0.8427,0.7974) (0.8347,0.8036) (0.8270,0.8101) (0.8196,0.8170) (0.8126,0.8243) (0.8059,0.8319) (0.7996,0.8398) (0.7936,0.8480) (0.7880,0.8566) (0.7828,0.8655) (0.7779,0.8746) (0.7734,0.8841) (0.7693,0.8939) (0.7656,0.9039) (0.7623,0.9142) (0.7593,0.9247) (0.7568,0.9356) (0.7546,0.9466) (0.7529,0.9578) (0.7516,0.9693) (0.7507,0.9810) (0.7502,0.9928) (0.7502,1.0048) (0.7506,1.0169) (0.7514,1.0291) (0.7527,1.0414) (0.7544,1.0538) (0.7566,1.0663) (0.7592,1.0787) (0.7623,1.0912) (0.7659,1.1036) (0.7699,1.1160) (0.7744,1.1283) (0.7795,1.1404) (0.7850,1.1524) (0.7909,1.1642) (0.7974,1.1757) (0.8044,1.1870) (0.8119,1.1980) (0.8198,1.2087) (0.8283,1.2189) (0.8372,1.2287) (0.8467,1.2381) (0.8566,1.2469) (0.8670,1.2552) (0.8778,1.2630) (0.8891,1.2700) (0.9009,1.2765) (0.9130,1.2822) (0.9256,1.2872) (0.9385,1.2914) (0.9518,1.2948) (0.9655,1.2973) (0.9794,1.2991) (0.9936,1.2999) (1.0080,1.2999) (1.0226,1.2989) (1.0373,1.2970) (1.0521,1.2943) (1.0670,1.2906) (1.0819,1.2860) (1.0967,1.2805) (1.1114,1.2741) (1.1259,1.2668) (1.1402,1.2588) (1.1542,1.2499) (1.1679,1.2403) (1.1812,1.2299) (1.1939,1.2189) (1.2062,1.2073) (1.2179,1.1950) (1.2290,1.1823) (1.2394,1.1691) (1.2491,1.1554) (1.2580,1.1414) (1.2662,1.1272) (1.2735,1.1127) (1.2799,1.0980) (1.2855,1.0832) (1.2902,1.0683) (1.2940,1.0534) (1.2968,1.0386) (1.2988,1.0238) (1.2998,1.0092) (1.3000,1.0000)};
\draw[->,ink,thick] (1.30,1)--(1.30,.91);
\draw[ink,thick] plot[smooth] coordinates {(1.7000,1.0000) (1.7000,0.9972) (1.6984,0.9643) (1.6942,0.9326) (1.6875,0.9021) (1.6784,0.8730) (1.6671,0.8453) (1.6537,0.8189) (1.6384,0.7939) (1.6213,0.7703) (1.6027,0.7480) (1.5826,0.7270) (1.5613,0.7073) (1.5388,0.6889) (1.5153,0.6716) (1.4911,0.6556) (1.4661,0.6407) (1.4405,0.6269) (1.4145,0.6142) (1.3882,0.6025) (1.3616,0.5917) (1.3348,0.5819) (1.3080,0.5730) (1.2813,0.5650) (1.2546,0.5577) (1.2280,0.5513) (1.2017,0.5457) (1.1756,0.5407) (1.1499,0.5365) (1.1245,0.5330) (1.0995,0.5302) (1.0749,0.5280) (1.0508,0.5264) (1.0271,0.5254) (1.0040,0.5250) (0.9813,0.5252) (0.9592,0.5259) (0.9377,0.5272) (0.9167,0.5291) (0.8962,0.5315) (0.8764,0.5344) (0.8571,0.5378) (0.8383,0.5417) (0.8202,0.5462) (0.8026,0.5512) (0.7856,0.5567) (0.7692,0.5626) (0.7533,0.5691) (0.7380,0.5761) (0.7233,0.5836) (0.7091,0.5916) (0.6955,0.6001) (0.6824,0.6092) (0.6699,0.6187) (0.6579,0.6288) (0.6464,0.6394) (0.6355,0.6505) (0.6251,0.6621) (0.6152,0.6743) (0.6059,0.6870) (0.5970,0.7003) (0.5887,0.7141) (0.5809,0.7285) (0.5735,0.7434) (0.5667,0.7589) (0.5604,0.7750) (0.5546,0.7916) (0.5493,0.8089) (0.5446,0.8266) (0.5403,0.8450) (0.5365,0.8639) (0.5333,0.8834) (0.5305,0.9035) (0.5284,0.9242) (0.5267,0.9454) (0.5256,0.9671) (0.5250,0.9894) (0.5251,1.0122) (0.5257,1.0355) (0.5269,1.0594) (0.5287,1.0836) (0.5311,1.1084) (0.5342,1.1335) (0.5380,1.1591) (0.5424,1.1849) (0.5476,1.2111) (0.5535,1.2375) (0.5602,1.2641) (0.5678,1.2908) (0.5761,1.3177) (0.5853,1.3444) (0.5955,1.3711) (0.6066,1.3977) (0.6186,1.4239) (0.6317,1.4498) (0.6459,1.4751) (0.6612,1.4999) (0.6777,1.5238) (0.6953,1.5470) (0.7142,1.5690) (0.7343,1.5900) (0.7558,1.6095) (0.7786,1.6276) (0.8027,1.6441) (0.8282,1.6587) (0.8551,1.6714) (0.8833,1.6819) (0.9129,1.6902) (0.9438,1.6960) (0.9759,1.6993) (1.0093,1.6999) (1.0437,1.6977) (1.0792,1.6927) (1.1154,1.6849) (1.1524,1.6741) (1.1900,1.6604) (1.2278,1.6438) (1.2658,1.6245) (1.3036,1.6024) (1.3411,1.5778) (1.3779,1.5508) (1.4138,1.5216) (1.4486,1.4904) (1.4820,1.4575) (1.5136,1.4231) (1.5434,1.3874) (1.5710,1.3508) (1.5962,1.3135) (1.6189,1.2757) (1.6390,1.2378) (1.6563,1.1999) (1.6708,1.1622) (1.6823,1.1251) (1.6910,1.0886) (1.6967,1.0529) (1.6996,1.0182) (1.7000,1.0000)};
\draw[->,ink,thick] (1.70,1)--(1.70,.91);
\draw[ink,thick] plot[smooth] coordinates {(2.3000,1.0000) (2.3000,0.9948) (2.2961,0.9347) (2.2857,0.8785) (2.2696,0.8262) (2.2481,0.7778) (2.2218,0.7330) (2.1913,0.6917) (2.1571,0.6538) (2.1198,0.6190) (2.0798,0.5872) (2.0375,0.5581) (1.9935,0.5316) (1.9480,0.5074) (1.9014,0.4853) (1.8541,0.4653) (1.8064,0.4471) (1.7583,0.4306) (1.7103,0.4157) (1.6625,0.4023) (1.6150,0.3901) (1.5680,0.3792) (1.5216,0.3694) (1.4759,0.3607) (1.4310,0.3529) (1.3870,0.3460) (1.3439,0.3400) (1.3019,0.3348) (1.2608,0.3303) (1.2208,0.3265) (1.1819,0.3234) (1.1441,0.3209) (1.1073,0.3189) (1.0717,0.3176) (1.0371,0.3167) (1.0036,0.3164) (0.9712,0.3166) (0.9399,0.3173) (0.9096,0.3184) (0.8804,0.3201) (0.8522,0.3221) (0.8249,0.3246) (0.7987,0.3276) (0.7734,0.3310) (0.7490,0.3348) (0.7255,0.3390) (0.7029,0.3437) (0.6812,0.3488) (0.6603,0.3544) (0.6403,0.3604) (0.6210,0.3668) (0.6025,0.3737) (0.5848,0.3811) (0.5678,0.3889) (0.5515,0.3972) (0.5359,0.4060) (0.5209,0.4153) (0.5066,0.4251) (0.4930,0.4355) (0.4799,0.4463) (0.4674,0.4578) (0.4556,0.4698) (0.4442,0.4823) (0.4335,0.4955) (0.4232,0.5093) (0.4135,0.5237) (0.4043,0.5388) (0.3956,0.5545) (0.3874,0.5709) (0.3797,0.5881) (0.3724,0.6060) (0.3656,0.6246) (0.3592,0.6440) (0.3533,0.6642) (0.3478,0.6853) (0.3428,0.7071) (0.3382,0.7299) (0.3340,0.7535) (0.3303,0.7781) (0.3270,0.8036) (0.3241,0.8300) (0.3217,0.8574) (0.3197,0.8858) (0.3182,0.9153) (0.3171,0.9458) (0.3165,0.9773) (0.3164,1.0099) (0.3169,1.0436) (0.3178,1.0783) (0.3193,1.1142) (0.3213,1.1511) (0.3239,1.1892) (0.3272,1.2283) (0.3311,1.2685) (0.3357,1.3098) (0.3411,1.3520) (0.3473,1.3953) (0.3543,1.4394) (0.3622,1.4845) (0.3712,1.5303) (0.3812,1.5768) (0.3923,1.6239) (0.4047,1.6715) (0.4184,1.7194) (0.4336,1.7674) (0.4504,1.8154) (0.4689,1.8631) (0.4893,1.9103) (0.5118,1.9567) (0.5364,2.0019) (0.5634,2.0457) (0.5930,2.0875) (0.6254,2.1271) (0.6608,2.1639) (0.6993,2.1974) (0.7412,2.2271) (0.7867,2.2525) (0.8358,2.2731) (0.8889,2.2882) (0.9458,2.2973) (1.0066,2.3000) (1.0714,2.2957) (1.1398,2.2841) (1.2118,2.2650) (1.2868,2.2380) (1.3645,2.2033) (1.4443,2.1610) (1.5252,2.1113) (1.6067,2.0547) (1.6876,1.9919) (1.7670,1.9235) (1.8439,1.8505) (1.9173,1.7739) (1.9860,1.6947) (2.0494,1.6139) (2.1066,1.5324) (2.1569,1.4514) (2.1999,1.3715) (2.2353,1.2936) (2.2629,1.2183) (2.2827,1.1460) (2.2950,1.0773) (2.2999,1.0122) (2.3000,1.0000)};
\draw[->,ink,thick] (2.30,1)--(2.30,.91);
\fill (1,1) circle (1.5pt);
\node[below] at (1,0) {$c/d=1$};
\node[left] at (0,1) {$a/b=1$};
\end{tikzpicture}

{\small Closed orbits around the centre $(1,1)$, for $a=b=c=d=1$.}
\end{center}


\subsection{Average populations over a cycle}
For a periodic positive orbit of period $T$,
\[
\frac{\dot y}y=c-dx\quad\Longrightarrow\quad
\frac1T\int_0^T x(t)\,dt
=\frac cd-\frac{\log y(T)-\log y(0)}{dT}=\frac cd.
\]
Similarly, integrating $\dot x/x=-a+by$ gives
\[
\overline x=\frac cd,\qquad\overline y=\frac ab.
\]
The time averages equal the equilibrium populations although the instantaneous
values oscillate. For the harvested model with positive coexistence, the
same calculation uses $a+h_1$ and $c-h_2$.

\section{The SIR epidemic model}
This is a mathematical compartment model under idealised assumptions,
not a fitted description of a particular current epidemic.
Let $S$ denote susceptible individuals, $I$ infectious individuals, and
$R$ removed individuals (recovered, isolated or dead).
Assume a closed, well-mixed population, no demographic turnover or return
to susceptibility, and transitions $S\to I\to R$.
The model is
\begin{equation}\label{eq:applied-sir}
\dot S=-\alpha SI,\qquad
\dot I=\alpha SI-\beta I,\qquad
\dot R=\beta I,\qquad\alpha,\beta>0.
\end{equation}
With population counts as variables, $\alpha$ has units
$(\text{population}\times\text{time})^{-1}$ and $\beta$ has units inverse time.
Some other conventions use an infection term divided by population size;
we consistently follow the supplied equations without that divisor.

Take $S(0)=S_0>0$, $I(0)=I_0>0$ and $R(0)=0$.
Adding equations gives $S+I+R=N=S_0+I_0$, a conserved total.
Thus $R=N-S-I$ and the physical $S$--$I$ phase region is the triangle
\[
S\geq0,\qquad I\geq0,\qquad S+I\leq N.
\]
The full line $I=0$ consists of equilibria. The $S$-nullclines are $S=0$
and $I=0$; the $I$-nullclines are $I=0$ and $S=\rho$, where
\[
\rho=\frac\beta\alpha.
\]
Inside the positive triangle, $\dot S<0$, so every trajectory moves left.
If $S>\rho$, $\dot I>0$; if $S<\rho$, $\dot I<0$.
The dimensionless initial growth threshold is $\alpha S_0/\beta=S_0/\rho$.

\subsection{An exact phase-plane relation}
For $S,I>0$,
\[
\frac{dI}{dS}=-1+\frac\rho S.
\]
Integrating and applying initial conditions gives
\begin{equation}\label{eq:applied-sir-phase}
I(S)=N-S+\rho\log\frac S{S_0}.
\end{equation}
This relates the two populations without explicitly solving for time.
% !TeX root = ../applied-mathematics1-companion.tex
% Dimensionless illustration N=1, S0=.99, I0=.01, rho=.3; not observed data.
\begin{center}
\begin{tikzpicture}[x=7.4cm,y=5cm,>=Stealth]
\draw[->] (0,0)--(1.06,0) node[right] {$S/N$};
\draw[->] (0,0)--(0,1.06) node[above] {$I/N$};
\draw[muted,dotted] (0,1)--(1,0);
\draw[muted,dashed] (.3,0)--(.3,.7);
\draw[very thick,ink,domain=0.0404097352:.99,samples=100]
 plot (\x,{1-\x+.3*ln(\x/.99)});
\draw[->,thick,ink] (.72,{1-.72+.3*ln(.72/.99)})
 --(.66,{1-.66+.3*ln(.66/.99)});
\draw[->,thick,ink] (.15,{1-.15+.3*ln(.15/.99)})
 --(.12,{1-.12+.3*ln(.12/.99)});
\fill (.99,.01) circle (1.5pt) node[above left,yshift=2pt] {initial};
\fill (.3,{.7+.3*ln(.3/.99)}) circle (1.5pt) node[above right] {peak};
\node[below] at (.3,0) {$\rho/N$};
\node[above right,muted] at (.68,.32) {$S+I=N$};
\node[below left] at (0,0) {$0$};
\end{tikzpicture}

{\small Illustrative SIR trajectory; arrows move towards decreasing $S$.}
\end{center}


\subsection{When does the infection peak?}
If $S_0>\rho$, $I$ initially rises. Its subsequent peak occurs when
$S=\rho$, giving
\[
I_{\max}=N-\rho+\rho\log\frac\rho{S_0}.
\]
If $S_0\leq\rho$, the trajectory starts on or to the left of this threshold,
so its maximum for $t\geq0$ is $I_0$, not the value of an unreachable
turning point obtained by extending the phase curve backwards.
The peak formula therefore needs the outbreak condition $S_0>\rho$.

\subsection{Removed population and final size}
Taking the ratio of the $S$ and $R$ equations gives
\[
\frac{dS}{dR}=-\frac S\rho
\quad\Longrightarrow\quad S(R)=S_0e^{-R/\rho}.
\]
Since $0\leq R\leq N$,
$S(t)\geq S_0e^{-N/\rho}>0$. Susceptibles are not exhausted completely;
infections stop growing when their reproduction rate falls below removals.
Substitution into $\dot R=\beta I$ gives the exact scalar equation
\begin{equation}\label{eq:applied-sir-R}
\dot R=\beta\bigl[N-R-S_0e^{-R/\rho}\bigr],\qquad R(0)=0.
\end{equation}
It is suitable for numerical integration but does not generally have a
simple elementary explicit solution.

\textit{Supplementary final-size relation.} Since $S$ decreases and $R$
increases while both remain bounded, they have limits $S_\infty,R_\infty$.
Then $I=N-S-R$ has a limit too. That limit must be zero, since
$\dot R=\beta I$ and a positive limiting $I$ would force unbounded $R$.
Therefore
\[
S_\infty=S_0e^{-(N-S_\infty)/\rho},\qquad R_\infty=N-S_\infty.
\]
This implicit equation determines the physical final state reached from
the given initial conditions; it is not an instruction to choose an
arbitrary algebraic root.

\section{An approximate explicit SIR solution}
The sources obtain a solvable approximation by assuming $R/\rho\ll1$.
Expanding the exponential in \eqref{eq:applied-sir-R} gives
\[
e^{-R/\rho}=1-\frac R\rho+\frac{R^2}{2\rho^2}
+O\bigl((R/\rho)^3\bigr).
\]
With $I_0=N-S_0$, define $q=S_0/\rho-1$ and $a=S_0/(2\rho^2)$.
The quadratic approximation is
\begin{equation}\label{eq:applied-sir-approx}
\dot R_{\mathrm{app}}=\beta(I_0+qR_{\mathrm{app}}-aR_{\mathrm{app}}^2),
\qquad R_{\mathrm{app}}(0)=0.
\end{equation}
Using a separate subscript emphasises that this is a different differential
equation, approximating the exact model only where the expansion is accurate.

\subsection{Solving the quadratic equation}
The source marks the detailed integration non-examinable. Retain it as an
optional worked derivation. Put
\[
\kappa=\sqrt{q^2+4aI_0},\qquad
b=\frac q{2a},\qquad A=\frac\kappa{2a}.
\]
Since $I_0>0$, $\kappa>|q|$ and hence $|q/\kappa|<1$.
Completing the square gives
\[
I_0+qR-aR^2=a\bigl[A^2-(R-b)^2\bigr].
\]
Let $z=(R-b)/A$. Then
\[
\dot z=\frac{\beta\kappa}{2}(1-z^2),\qquad
\operatorname{artanh}z=\frac{\beta\kappa t}{2}+C.
\]
At $t=0$, $z=-q/\kappa$, so set
$\phi=\operatorname{artanh}(q/\kappa)$ to obtain
\begin{equation}\label{eq:applied-sir-tanh}
R_{\mathrm{app}}(t)=\frac{\rho^2}{S_0}
\left[q+\kappa\tanh\left(\frac{\beta\kappa t}{2}-\phi\right)\right].
\end{equation}
Here $\operatorname{artanh}$ is the inverse hyperbolic tangent, not the
reciprocal of $\tanh$. Differentiating the approximate solution gives
\[
\dot R_{\mathrm{app}}(t)=\frac{\beta\kappa^2\rho^2}{2S_0}
\operatorname{sech}^2\left(\frac{\beta\kappa t}{2}-\phi\right).
\]
The initial condition is a useful check: $\tanh(-\phi)=-q/\kappa$, so
$R_{\mathrm{app}}(0)=0$.

\begin{remark}[Keep the approximation consistent]
For the exact model, $S=S_0e^{-R/\rho}$ and $I=N-R-S$.
Inserting $R_{\mathrm{app}}$ into these expressions gives approximate
reconstructions, not an exact SIR solution. Alternatively
$I_{\mathrm{app}}=\dot R_{\mathrm{app}}/\beta$ follows from the quadratic
truncation. These two reconstructions agree only to the retained order:
they should not be described as identical exact functions.
A large numerical $\rho$ is not sufficient by itself; the dimensionless
ratio $R/\rho$ must remain small over the time range of interest.
\end{remark}

\subsection{What the source's historical illustrations show}
For the supplied boarding-school illustration, the stated inputs are
$N=760$, $I_0=1$, $S_0=759$ and a peak near $300$ infections.
Using the exact phase-plane peak formula gives an equation for $\rho$:
\[
760-\rho+\rho\log(\rho/759)=300.
\]
On $0<\rho<759$, numerical solution gives roughly $\rho\approx195.0$.
The source rounds this to about $200$. The peak height determines this
parameter ratio, but not the time scale: one needs timing information to
identify $\beta$, then $\alpha=\beta/\rho$.

For the source's rounded choice $\rho=200$, $q\approx2.795$,
$\kappa\approx2.802$, and $q/\kappa\approx0.998$.
Thus $\phi$ is around $3.4$, not $\operatorname{artanh}(0.014)$ as printed.
Also the total number removed can be much larger than $200$, so the
small-$R/\rho$ approximation cannot be assumed valid over the whole outbreak.
These are model illustrations from the handout; the underlying data are
not supplied as a table, and no data fit has been reproduced here.

The Bombay plague illustration similarly compares a curve proportional
to $\dot R$, rather than cumulative $R$, with deaths per unit time.
Interpreting deaths as a fixed fraction of removals requires an extra
observation assumption; the compartment equations alone do not provide
that conversion. The source's plotted observations and fitted parameters
are not reconstructed as new evidence.

\section{A second-order equation as a planar system}
For
\[
\ddot x+3\dot x+2\sin x=0,
\]
set $y=\dot x$. Then
\[
\dot x=y,\qquad\dot y=-3y-2\sin x.
\]
The nullclines are $y=0$ and $y=-(2/3)\sin x$, and the equilibria are
$(n\pi,0)$ for integers $n$. At such a point,
\[
J=\begin{pmatrix}0&1\\-2\cos(n\pi)&-3\end{pmatrix},\qquad
\lambda^2+3\lambda+2\cos(n\pi)=0.
\]
For even $n$, the eigenvalues are $-1,-2$, giving local attraction.
For odd $n$, they are $(-3\pm\sqrt{17})/2$, with opposite signs,
giving an unstable saddle.
This conversion lets the same equilibrium and Jacobian methods analyse
second-order motion, connecting this chapter with mechanics.

\studyguide{For a continuous model: find zero rates, not intersections
with $y=x$. Use derivative signs in one dimension and Jacobian eigenvalue
real parts in two dimensions. If the linear test has zero real parts,
look for nonlinear information such as a first integral. State any domain
or small-parameter assumptions before using an explicit approximation.}
